\section{Discussion}\label{sec:discussion} 
Observation enters the problem via (i) nonresolution, which leaves an unknown amount of terminal bad mass, and (ii) the resolved sub-sequence, which defines a known information quantity. The profile representation splits the information into the part absorbed by resolved observation times $\tau_\Pi$, and the residual which then bounds terminal resolution. Any coarsening decreases the former information and thus allows a higher residual to terminal bad events, increasing the terminal event pool compatible with a given total information level.

Knowing the full curve $\rho_\Pi^*(\beta)$ allows immediate assessment of at what levels of information budget the conclusion of a $\beta$-admissible risk level is forced (or excluded) by an observed distribution. Coarsening the timing information therefore moves this curve according to the additional information per the profile representation, as seen in the examples of decisions on the population risk and certification of the absence of power under endpoint non-certifiability, respectively.

Certifying against timing distributions with finitely many timing levels inherently incurs information cost: the observed data allowed much shorter certification horizons under the timing-distribution maximiser (with 2 timing bands), compared to the null-equivalent distribution supporting endpoint-only decisions. Further levels of timing refinement could lead to tighter intervals on the population quantity at the risk of wider confidence regions, increased relaxation parameter and frequency of failure of the numerics. We therefore recommend choosing an information budget to match anticipated timings a priori, and against post-hoc best-case over the family of non-adjusted refinements of timing.

All results above assume faithful classification labelling and event observations up to a known maximum age of interest, a static multi-class distribution of events, and satisfy a condition on the information bound wrt. a hypothetical continuous stream of interest. The example laws illustrate our argument, but we have made no attempt to estimate how widely this scenario applies in practice, nor obtain universal bounds for admissibility of $\rho$. Improving the observable confidence set construction and/or tightening the projections, and problem-informed specification of the information assumption are all avenues for future work; none of which we can comment on beyond re-visiting these assumptions. The projection construction is rigorous w.r.t. its underlying identities. For the accelerated construction of confidence intervals, we additionally assume valid comparisons under the floating point precision and truncations of filters in use, which our record of runs does not allow us to uniformly verify. We take the suite of simulation studies to indicate correct behaviour, without certifying the arithmetic identities underlying our claim about any individual simulation.

One practical scenario where the framework would be useful is tracking a batch process which admits early/late/non-terminal information, relative to an upper age cut-off. This requires maintaining the timing stratification of each observed event, assurances about labelling accuracy of events, and communication of an admissible class of values for the cumulative information bound given a decision on maximum $\beta$. Certification is deferred until reaching that known age, but thereafter can be considered permanently valid. We do not assess application of the tool to any running example.

\section{Conclusion}\label{sec:conclusion} 
We have shown that under a bounded mutual information from a categorised information disclosure to latent resolution, an upper and lower bound on maximum terminal hazard is determined by the single value of the corresponding information-completion. More information along the disclosed sequence strictly lowers the minimum information needed for compatibility and translates directly into the boundary of information-compatibility for each decision threshold. Using a joint confidence sequence on the disclosed distribution and mapping via this information-boundary enables an any-time certificate, subject to valid enclosing properties of the numerical estimates. We showed empirically that observing intermediate disclosures permits decisions on the endpoint distribution inaccessible to a population-based analysis alone, and achieves much higher confidence at significantly reduced time under non-degenerate disclosed information. Our negative control demonstrates the necessity of having sufficient timing information, against costs in confidence estimation, relaxation of the target, and incomplete computation.
